% vim: ai sw=2 spell spelllang=cs

\iffimuni\else

\begin{frame}{Práce s pamětí}

  \heading{LDD3 kap. 8} (zastaralá část o bootmem)\\
  \heading{Understanding the Linux Virtual Memory Manager} (zastaralá)

  \bigskip

  \begin{itemize}
  \item Organizace paměti
%  \item Typy paměti
  \item Alokace paměti
  \end{itemize}

\end{frame}

\fi

\iffimuni
\section{Organizace paměti}
\frame{\sectionpage}
\fi

\begin{frame}
  \frametitle{Organizace paměti}

  \begin{center}
  \heading{2 oddělené světy} \\
  (tentokrát z pohledu HW)

  \pause
  \bigskip

  tj.\\
  \heading{2 různé adresové prostory}
  \end{center}

  \pause
  \medskip

  \heading{Fyzický prostor}
  \begin{itemize}
  \item Rozložení určuje BIOS (popř. ekvivalenty -- (U)EFI, Qemu)
  \item x86: \texttt{dmesg|grep BIOS-e820}
  \end{itemize}

  \pause

  \heading{Virtuální prostor (paměť)}
  \begin{itemize}
  \item Rozložení si určuje OS
  \begin{itemize}
  \item Různé na každé architektuře
  \end{itemize}
  \item Dané mapování mezi FP a VP
  \item Nutná podpora HW (MMU), jinak mapování 1:1
  \end{itemize}

\end{frame}

\begin{frame}{Organizace paměti -- graficky}

  \begin{center}
    % vim: ai sw=2 spell spelllang=cs
%
\tikzstyle{spcwid} = [minimum width=4cm]
\tikzstyle{spc} = [draw,thick,spcwid,minimum height=.85\textheight]
\tikzstyle{box} = [draw,fill=white]
\tikzstyle{empty} = [inner sep=0pt,minimum width=0pt]
\tikzstyle{line} = [-,semithick]
\pgfdeclarelayer{ram}
\pgfdeclarelayer{space}
\pgfsetlayers{ram,main,space}

\begin{tikzpicture}
  \begin{pgfonlayer}{space}
    \node [spc] (PHYS) {};
  \end{pgfonlayer}

  \node [rotate=90,left=0mm of PHYS,anchor=south] {Physical Space};
  \begin{pgfonlayer}{ram}
    \node [draw,fill=blue!15,spcwid,minimum height=5.4cm,below=0mm of PHYS.north] (RAM) {RAM};
  \end{pgfonlayer}

  \node [box,below right=0mm of PHYS.north west] {BIOS};
  \node [box,below right=2.8cm and 2mm of PHYS.north west] {ACPI};
  \node [box,below right=3.5cm and 1.5cm of PHYS.north west] {Timer};
  \node [box,spcwid,below=4.5cm of PHYS.north] {PCI Devices};

  \path (RAM.south) -- node {Empty} (PHYS.south);

  \begin{pgfonlayer}{space}
    \node [spc,right=18mm of PHYS] (VIRT) {};
  \end{pgfonlayer}
  \node [rotate=270,right=0mm of VIRT,anchor=south] {Virtual Space};

  \node [box,spcwid,minimum height=.47\textheight,align=center,below=0mm of VIRT.north] (US)
    {Userspace \\ (Differs per Process)};
  \begin{pgfonlayer}{ram}
    \node [box,fill=blue!15,spcwid,minimum height=3.4em,below=-\pgflinewidth of US.south,align=center] (PHYS2) {Physical \\ Memory (RAM, \dots)};
  \end{pgfonlayer}
  \node [box,below right=-\pgflinewidth and 0pt of US.south west] (HOLE) {Hole};
  \node [spcwid,above=-\pgflinewidth of VIRT.south] (MODS) {Modules};
  \node [box,spcwid,above=-\pgflinewidth of MODS.north] (KERN) {Kernel};
  \node [box,minimum width=1.5cm,above left=-\pgflinewidth and 0pt of KERN.north east] {Hole};
  \path (PHYS2.south) -- node {\dots} (KERN.north);

  \node [empty,below=1cm of RAM.north east] (CONNL1) {};
  \node [empty,below=2mm of CONNL1] (CONNL2) {};

  \node [empty,left=1cm of CONNL1] (RECTL) {}
    edge[line,densely dashed] (CONNL1);
  \node [empty,below=8mm of RECTL] (RECBL) {}
    edge[line,densely dashed] (RECTL);
  \node [empty,right=8mm of RECBL] (RECBR) {}
    edge[line,densely dashed] (RECBL);
  \node [empty,left=2mm of CONNL2] (RECTR) {}
    edge[line,densely dashed] (CONNL2)
    edge[line,densely dashed] (RECBR);

  \node [empty,below=1cm of US.north west] (CONNR1) {}
    edge[line] (CONNL1);
  \node [empty,below=2mm of CONNR1] (CONNR2) {}
    edge[line] (CONNL2);

  \draw (RAM.north east) -- (HOLE.south west);
  \draw (RAM.south east) -- (PHYS2.south west);

  \node [empty,above=1cm of PHYS.south east] (ARRL) {};
  \node [empty,above=1cm of VIRT.south west] {}
    edge[<->,semithick,>=stealth']
      node[above] {\footnotesize MMU}
      node[below] {\footnotesize Mapping} (ARRL);
\end{tikzpicture}

  \end{center}

\end{frame}

\iffimuni
\subsection{Fyzický prostor}
\fi

\begin{frame}
  \frametitle{Fyzický prostor}

  \begin{itemize}
    \item To, co je za MMU
    \item RAM, časovače, řadiče přerušení, ACPI tabulky, zařízení

    \pause

    \item Co kde je $\rightarrow$ otázka na BIOS/(U)EFI/\dots

    \pause

    \item Jak který prostor vypadá $\rightarrow$ specifikace bridge, ACPI, PCI,
      \ldots

    \pause

    \item I hierarchický prostor
    \begin{itemize}
    \item Např.\ rozsah \code{100}--\code{1ff} $\rightarrow$ ChipXY
      \begin{itemize}
      \item ChipXY směruje dále:
      \item Adresa \code{100}--\code{18f} $\rightarrow$ řadič klávesnice
      \item Adresa \code{190}--\code{1ff} $\rightarrow$ časovač
      \end{itemize}
    \end{itemize}
  \end{itemize}

  \pause

  \heading{RAM}
  \begin{itemize}
    \item Reprezentována rozsahy lineárních (fyzických) adres
    \begin{itemize}
      \item Rozsahy kvůli vloženým prostorům zařízení (e820)
    \end{itemize}
    \item Cache (některá zařízení také)
  \end{itemize}

\end{frame}

\begin{frame}
  \frametitle{Úkol}
  
  \heading{Práce s fyzickou adresou -- čtení signatury a délky ACPI
    tabulky z fyzické adresy}
  \begin{enumerate}
    \item Najděte fyzickou adresu nějaké ACPI tabulky v \cmd{dmesg}
    \begin{itemize}
      \item Řádky formátu \code{ACPI: XSDT 00000000af5b0100 00064}
      \item Jeden si zvolte (\emph{jinou než} \code{RSDP})
      \item Zapamatujte si adresu (\code{0x00000000af5b0100})
      \item Datový typ pro fyzickou adresu je \code{phys_addr_t}
    \end{itemize}

    \item Získejte virtuální adresu (\code{u32 *virt = ioremap_cache(phys, 8)})
    \item Vypište virtuální adresu (\code{\%p} a porovnejte hodnotu s fyzickou)
    \item Vypište signaturu (4\,B textu -- \code{virt[0]} a \code{\%.4s})
    \item Vypište délku tabulky (4\,B -- \code{virt[1]} a \code{\%u})
    \item Odmapujte (\code{iounmap(virt)})
    \item Porovnejte délku s výpisem (tj. \code{0x00064} nahoře)
    \item Můžete si přemapovat a vypsat celou tabulku a porovnat položky se
      specifikací ACPI
  \end{enumerate}

  \bigskip

  Pozn.: na starších jádrech se výpis v \cmd{dmesg} liší

\end{frame}

\iffimuni
\subsection{Virtuální prostor}
\fi

\begin{frame}
  \frametitle{Virtuální prostor}

  \begin{itemize}
    \item To, co vidí CPU (a překládá MMU)
    \begin{itemize}
      \item I 1:1 -- start systému, noMMU, \ldots
    \end{itemize}

    \pause

    \item OS ,,rozdělí'' FP na stránky
    \begin{itemize}
      \item Index stránky je tzv. \emph{page frame number} (pfn)
    \end{itemize}

    \pause

    \item OS namapuje stránky do VP
    \begin{itemize}
      \item A může přidat ,,imaginarní'' stránky např.\ na disku (swap)
      \begin{itemize}
	\item Programy vidí více ,,paměti''
      \end{itemize}

      \item Toto mapování čte MMU (ví odkud)
    \end{itemize}

    \pause

    \item Každá stránka má svoji \code{struct page}
    \begin{itemize}
      \item Informace o stránce (nikdy nevyswapovat, špinavá, počet uživatelů)
    \end{itemize}
  \end{itemize}

\end{frame}

\begin{frame}
  \frametitle{Grafické shrnutí}

  \begin{center}
  \resizebox{!}{.85\textheight}{
    \pgfdeclarelayer{background}
    \pgfsetlayers{background,main}

    \tikzstyle{ent} = [rectangle, thick, draw, inner sep=12pt,fill=white]
%      fill=orange!25]
    \tikzstyle{mmu} = [ent, rounded corners=8pt, minimum width=2.2cm]
    \tikzstyle{arr} = [->, >=stealth', thick]
    \tikzstyle{note} = [font=\footnotesize]
    \tikzstyle{empty} = [inner sep=0pt, minimum size=0pt]
    \tikzstyle{fillline} = [thick, dotted]

    \begin{tikzpicture}[node distance=8mm]
        \node [ent, minimum width=6cm] (MEM) { RAM };
        \node [note] (PHYS) [below=1.5mm of MEM] { Physical addresses };

	\node [empty] (MEML) [below left=0mm and -1cm of MEM] {};
	\node [empty] (MEMR) [below right=0mm and -1cm of MEM] {};

        \begin{uncoverenv}<2->
        \node [mmu] (IOMMU) [below=of MEML] { IOMMU }
	  edge[arr] node {} (MEML);
        \node [ent] (DEV) [below=of IOMMU, align=center] {
		\begin{tabular}{c@{}c@{}}
			Device\\[-1.5mm]
			\tiny{(DMA)}
		\end{tabular}
	  }
	  edge[arr] node [note, right, align=left] {Device \\ addresses} (IOMMU);
        \end{uncoverenv}

        \node [mmu] (MMU) [below=of MEMR] { MMU }
	  edge[arr] node {} (MEMR);
        \node [ent] (CPU) [below=of MMU, align=center] {
		\begin{tabular}{c@{}c@{}}
			CPU\\[-1.5mm]
			\tiny{(program)}
		\end{tabular}
	  }
	  edge[arr] node [left, note, align=right] {Virtual \\ addresses} (MMU);

	\node[empty] (CENR) [right=3mm of MMU] {};
	\node[empty] (CENL) [left=3mm of IOMMU] {};
	\node[empty] (TL) [above=2.8cm of CENL] {};
	\node[empty] (TR) [above=2.8cm of CENR] {};
	\node[empty] (CEN) [right=3.3cm of CENL] {};
	\node[empty] (BL) [below=3.1cm of CENL] {};
	\node[empty] (BR) [below=3.1cm of CENR] {};
	\node[empty] (BOT) [right=3.3cm of BL] {};

	\begin{pgfonlayer}{background}
	\fill[color=green!25] (TL) rectangle (CENR);
	\fill<2->[color=blue!25] (CENL) rectangle (BOT);
	\fill[color=red!25] (CEN) rectangle (BR);
	\draw[fillline,color=green] (CENL) -- (TL) -- (TR) -- (CENR);
	\draw<2->[fillline,color=blue] (CENL) -- (CEN) -- (BOT) -- (BL) -- (CENL);
	\draw[fillline,color=red] (CEN) -- (CENR) -- (BR) -- (BOT);
	\end{pgfonlayer}
    \end{tikzpicture}
  }~~~~~~~~~~~~~
  \end{center}

\end{frame}

\begin{frame}
  \frametitle{Operace s adresami}

  \begin{itemize}
    \item \header{linux/mm.h}, \header{linux/io.h}
    \item Např. \code{void *virt = __get_free_page}
  \end{itemize}

  \heading{Fyzická adresa}
  \begin{itemize}
    \item \code{phys_addr_t phys = virt_to_phys(virt)}
    \item \code{void *virt = phys_to_virt(phys)}
  \end{itemize}

  \heading{Struktura \code{page}}
  \begin{itemize}
    \item \code{struct page *page = virt_to_page(virt)}
    \item \code{void *virt = page_to_virt(page)}
  \end{itemize}

  \heading{PFN} (Page Frame Number)
  \begin{itemize}
    \item \code{unsigned long pfn = page_to_pfn(virt_to_page(virt))}
    \item \code{void *virt = pfn_to_virt(pfn)}
    \item \code{unsigned long pfn = phys >> PAGE_SHIFT}
    \item \code{phys_addr_t phys = pfn << PAGE_SHIFT}
  \end{itemize}

\end{frame}

\begin{frame}
  \frametitle{Úkol}
  
  \heading{Práce s virtuální adresou}
  \begin{enumerate}
  \item \header{linux/mm.h}
  \item Naalokujte stránku (\code{virt = __get_free_page(GFP_KERNEL)})
  \item Nakopírujte do ní řetězec
  \item Zjistěte fyzickou adresu (\code{phys = virt_to_phys(virt)})
  \item Zjistěte adresu {struct page} (\code{page = virt_to_page(virt)})
  \item Zarezervujte stránku (\code{SetPageReserved(page)})
  \item Přemapujte fyzickou stránku (\code{map = ioremap(phys, ...)})
  \begin{itemize}
    \item Varování v \cmd{dmesg} ignorujte
  \end{itemize}
  \item Vypište \code{virt}, \code{phys}, \code{page}, \code{map} a
    \code{page_to_pfn(page)}
  \item Vypište obsah (\code{\%s}) \code{virt} a \code{map}
  \item Změňte obsah \code{map} a vypište obsah \code{virt}
  \item \code{iounmap}, \code{ClearPageReserved} a \code{free_page}
  \item Zkonzultujte s \doc{Documentation/x86/x86\_64/mm.txt}
  \end{enumerate}
\end{frame}

\iffimuni
\section{Alokace paměti}
\frame{\sectionpage}
\fi

\begin{frame}
  \frametitle{Alokace paměti}

  \heading{Základní typy alokací}
  \begin{itemize}[<+->]
    \item Malé bloky až souvisle po stránkách (\code{kmalloc})
    \item Souvisle po stránkách (\code{__get_free_pages})
    \item Nesouvisle po stránkách a namapovat (\code{vmalloc})
  \end{itemize}

  \pause
  \medskip

  \heading{GFP\_*}
  \begin{itemize}
    \item Parametr pro alokátory
    \item \code{GFP_ATOMIC} -- nespi (např.\ uvnitř spinlocků)
    \item \code{GFP_KERNEL} -- obyčejné alokace všude jinde
    \item a další: \code{GFP_NOFS}, \code{GFP_NOIO}, \dots
  \end{itemize}

\end{frame}

\iffimuni
\subsection{SLAB alokátor}
\fi

\begin{frame}
  \frametitle{\code{kmalloc}}

  \heading{SLAB alokátor}
  \begin{itemize}
    \item Vnitřně alokuje po stránkách
    \begin{itemize}
      \item Stránka má velikost \code{PAGE_SIZE} (x86: 4K)
    \end{itemize}
    \item Rozděluje stránky na menší kusy
    \begin{itemize}
    \item Výhýbá se fragmentaci paměti
    \end{itemize}
    \item \header{linux/slab.h}
    \item \code{kmalloc}, \code{kzalloc} (= \code{kmalloc} +
      \code{memset}), \code{kfree}
  \end{itemize}

\end{frame}

\begin{frame}
  \frametitle{Úkol}

  \heading{Alokace pomocí \code{kmalloc}}

  \begin{enumerate}
    \item Alokujte 32 stránek (jednou alokací)
    \item Vypište jejich pfn
    \begin{itemize}
      \item Skákejte po \code{PAGE_SIZE}
      \item Použijte \code{virt_to_page} a \code{page_to_pfn}
    \end{itemize}
  \end{enumerate}

\end{frame}

\iffimuni
\subsection{Stránkový alokátor}
\fi

\begin{frame}
  \frametitle{\code{__get_free_pages}}

  \heading{Stránkový alokátor}
  \begin{itemize}
  \item Musí najít souvislý blok volných stránek
  \begin{itemize}
    \item Může být problém najít větší bloky ($\geq$ 4 stránky) kvůli
      fragmentaci
  \end{itemize}

  \item Rychlejší než \code{kmalloc}

  \item Používá tzv.\ \emph{buddy} systém
  \begin{itemize}
    \item Snižuje fragmentaci
    \item Podrobnosti viz wiki \iffimuni nebo Překladače \fi
  \end{itemize}

  \pause

  \item \header{linux/mm.h}
  \item Velikost alokace se udává řádem, \code{order} ($2^{order}$ = počet
    stránek)
  \begin{itemize}
    \item \code{get_order}, \code{MAX_ORDER} (na x86: 11, tj.\
      $2^{11} \cdot 4\text{k} = 8\text{M}$)
  \end{itemize}

  \item \code{__get_free_pages}, \code{free_pages}
  \item \code{__get_free_page}, \code{free_page} (zkratky pro řád \code{0})
  \end{itemize}

  \iffimuni
  \bigskip
  \pause

  \heading{Úkol:} alokace 32 stránek a výpis pfn
  \fi

\end{frame}

\iffimuni
\subsection{Virtuální alokátor}
\fi

\begin{frame}
  \frametitle{\code{vmalloc}}

  \heading{Virtuální alokátor}
  \begin{itemize}
  \item Použitím podobný \code{kmalloc}
  \item \pozor: neatomický (vnitřně používá \code{GFP_KERNEL})
  \item Maximální velikost alokace daná omezením architektury
  \begin{itemize}
    \item x86\_32: 128\,MiB (lze zvýšit parametrem, ne o moc)
    \item x86\_64: 30\,TiB
  \end{itemize}

  \pause

  \item Alokuje po stránkách a mapuje je souvisle
  \begin{itemize}
    \item Fyzické stránky jsou umístěné různě po RAM
    \item Dražší
  \end{itemize}

  \pause

  \item \header{linux/vmalloc.h}
  \item \code{vmalloc}, \code{vfree}
  \item Odlišná práce s adresami (\emph{POZOR})
  \begin{itemize}
  \item \code{vmalloc_to_page}, \code{vmalloc_to_pfn}
  \end{itemize}
  \end{itemize}

  \bigskip
  \pause

  \heading{Úkol:} alokace 32 stránek a výpis pfn

\end{frame}

\begin{frame}
  \frametitle{Poznámky na závěr}

  \begin{itemize}
  \item Vše, co může číst uživatel se musí nejprve vymazat
  \begin{itemize}
    \item Pomocníci: \code{kzalloc}, \code{get_zeroed_page}, \code{vzalloc}
    \item Nejlépe vymazat všechno, co alokuji a není kritické
  \end{itemize}
  \item Používat správné \code{GFP_*}
  \begin{itemize}
    \item \code{GFP_ATOMIC} funguje všude, ale ubírá vzácné prostředky
  \end{itemize}
  \end{itemize}

\end{frame}

\begin{frame}
  \frametitle{Úkol}

  \heading{Alokace paměti}

  \begin{enumerate}
  \item Zjistěte, kolikrát za sebou lze naalokovat pamět:
  \begin{itemize}
  \item řádu 10 jako \code{GFP_ATOMIC}
  \item řádu 10 jako \code{GFP_ATOMIC} (znovu po 5 vteřinách)
  \item řádu 10 jako \code{GFP_KERNEL}
  \item řádu 10 jako \code{GFP_KERNEL} po uvolnění předešlé paměti
  \item Pozn.: $2^{10} \cdot 4096 = 4\,\text{M}$, jak velké pole na ukazatele?
  \end{itemize}

  \item Zkuste ve smyčce alokovat paměť s řádem $\leq$
    \code{PAGE_ALLOC_COSTLY_ORDER}, co se stane?
    \pause
    Přijde OOM killer\dots
  \end{enumerate}

  \bigskip

  Inspirace: pb173/05

\end{frame}
